\documentclass[11pt,a4paper]{article}
\usepackage[utf8]{inputenc}
\usepackage[T1]{fontenc}
\usepackage{amsmath,amssymb}
\usepackage{siunitx}
\usepackage{booktabs}
\usepackage{geometry}
\usepackage{hyperref}
\geometry{margin=2.5cm}

\title{Single-Phase Power Transformer Sizing\\
\large A Practical Methodology for High-Frequency, High-Power Designs}
\author{}
\date{}

\begin{document}
\maketitle

\section{Scope}
This document summarizes the engineering steps required to dimension a single-phase
power transformer for switch-mode converter applications (e.g.\ PSFB, DAB, isolated
DC/DC stages), covering:
\begin{itemize}
    \item turns ratio determination from the voltage transfer function,
    \item core material selection as a function of operating frequency,
    \item core sizing via the area-product method,
    \item winding design considerations (skin/proximity effects, litz wire vs.\ foil),
    \item a worked numerical example at \SI{800}{V}/\SI{800}{V}, \SI{10}{kHz}, \SI{250}{kW}.
\end{itemize}

\section{Turns Ratio}
For an ideal isolated transformer driven by a quasi-square or square voltage waveform
(as typical of PSFB, DAB, or any hard/soft-switched bridge), the volt-second balance on
each winding gives Faraday's law in its waveform-corrected form:
\begin{equation}
    V = K_f \, N \, A_c \, B_{m} \, f
    \label{eq:faraday}
\end{equation}
where:
\begin{itemize}
    \item $V$ = RMS (or peak, depending on $K_f$ convention) voltage applied to the winding [V],
    \item $N$ = number of turns,
    \item $A_c$ = core cross-sectional area [m$^2$],
    \item $B_m$ = peak flux density [T],
    \item $f$ = switching/operating frequency [Hz],
    \item $K_f$ = waveform factor: $K_f = 4$ for a symmetric square wave, $K_f = 4.44$ for a pure sinusoid, intermediate values for trapezoidal/quasi-square waveforms with duty cycle $D<0.5$ (e.g.\ $K_f \approx 4D$ for a unipolar square wave of duty $D$).
\end{itemize}

The turns ratio is fixed by the required voltage transformation:
\begin{equation}
    n = \frac{N_1}{N_2} = \frac{V_1}{V_2}
\end{equation}
For converters requiring voltage step change under varying duty cycle (e.g.\ PSFB with
output inductor), $n$ is chosen at the nominal/maximum duty point, with margin for
line and load regulation.

\section{Core Material Selection vs.\ Frequency}
The dominant constraint at a given frequency is core loss, which is approximately:
\begin{equation}
    P_{core} \approx k_c \, f^{\alpha} \, B_m^{\beta} \, V_{core}
\end{equation}
(Steinmetz equation), where $\alpha,\beta$ and $k_c$ are material-dependent. The choice
of material is therefore frequency-driven:

\begin{center}
\begin{tabular}{@{}lllp{5.5cm}@{}}
\toprule
Material & Typical $f$ range & $B_{sat}$ [T] & Notes \\
\midrule
Grain-oriented Si-Fe (laminated) & DC--\SI{1}{kHz} & 1.6--2.0 & High saturation, but eddy-current losses scale with $f^2 t^2$; lamination thickness limits practical use above $\sim$\SI{1}{kHz}. \\
Amorphous (Fe-based) & \SI{1}{}--\SI{10}{kHz} & 1.5--1.6 & Lower loss than Si-Fe at the same $f$, good compromise for medium-frequency, high-power transformers. \\
Nanocrystalline (e.g.\ Vitroperm, Finemet) & \SI{1}{}--\SI{150}{kHz} & 1.2--1.5 & High $B_{sat}$ retained at MF, low specific losses; preferred for high-power MF transformers (DAB, solid-state transformers, traction). \\
MnZn ferrite & \SI{20}{kHz}--\SI{500}{kHz} & 0.35--0.5 & Low $B_{sat}$ limits power density at low kHz; becomes attractive above $\sim$\SI{50}{kHz} where its losses remain low and nanocrystalline/amorphous losses rise. \\
NiZn ferrite & $>$\SI{500}{kHz} & 0.3--0.4 & High-frequency, low-power applications; not relevant for high-power MF transformers. \\
\bottomrule
\end{tabular}
\end{center}

\textbf{Rule of thumb:} for a given frequency, select the material that allows the
highest usable $B_m$ (limited by core loss density, typically targeting
\SIrange{0.5}{1.5}{W/cm^3} depending on cooling) while satisfying thermal constraints.
At \SI{10}{kHz} and high power, nanocrystalline or amorphous tape-wound cores
typically outperform ferrite because their higher $B_{sat}$ allows a smaller, lighter
core for the same volt-seconds, even though their specific core losses per unit volume
at \SI{10}{kHz} are higher than ferrite's — the reduction in required volume usually
wins for power levels in the hundreds of kW.

\section{Core Sizing: Area-Product Method}
The area product $A_p = A_c \cdot A_w$ (core area times window area) relates the
apparent power transferred to the achievable current density and flux density,
and is the standard first-pass sizing tool (McLyman):
\begin{equation}
    A_p \;[\text{cm}^4] = \frac{P_t \cdot 10^4}{K_f \, K_u \, B_m \, f \, J}
    \label{eq:ap}
\end{equation}
where:
\begin{itemize}
    \item $P_t$ = total apparent power handled by the transformer, $P_t \approx P_{in}+P_{out} \approx 2P$ for an isolated transformer carrying power $P$ on each side,
    \item $K_u$ = window utilization factor ($\approx 0.2$--$0.3$ for round/litz wire with insulation and creepage, up to $\approx 0.4$--$0.5$ for foil/busbar windings),
    \item $B_m$ = peak flux density [T] (material- and loss-limited, see above),
    \item $f$ = frequency [Hz],
    \item $J$ = RMS current density [A/cm$^2$] (cooling-dependent: $\sim$200--350 for natural convection, up to 500+ with forced air/liquid cooling).
\end{itemize}
Once $A_p$ is known, a core (or stack of cores) is selected from a manufacturer
catalogue with matching or slightly larger $A_p$, giving $A_c$ and $A_w$ directly.
Equation~\eqref{eq:faraday} then yields the turns count for each winding.

\section{Winding Design}
At MF/HF, conduction losses are dominated by skin and proximity effects rather
than DC resistance:
\begin{equation}
    \delta = \sqrt{\frac{\rho}{\pi f \mu_0 \mu_r}}
\end{equation}
For copper, $\delta \approx 7.5/\sqrt{f\,[\text{Hz}]}$ cm. At \SI{10}{kHz},
$\delta \approx \SI{0.66}{mm}$, meaning conductor thickness should not exceed
roughly $2\delta \approx \SI{1.3}{mm}$ in any single layer.

\begin{itemize}
    \item \textbf{Litz wire}: appropriate for moderate current windings where many
    strands keep each strand below $\delta$; strand count and bundle geometry must
    be checked for proximity losses in multi-layer windings.
    \item \textbf{Copper foil}: preferred at very high current, low turn-count
    windings (as in the example below); foil thickness $\le 2\delta$, interleaved
    with thin insulation, often the only practical option above a few hundred
    amperes RMS per turn.
    \item \textbf{Interleaving}: primary/secondary interleaving reduces the MMF
    excursion across the window and is essential to control proximity losses at
    MF in high-current designs.
\end{itemize}

\section{Insulation and Isolation}
Creepage and clearance distances are dictated by the working voltage, altitude,
pollution degree, and required isolation class (basic/supplementary/reinforced)
per IEC~60664-1 or the relevant converter standard. Tape-wound cores with bobbin
or potted construction are common; for busbar/foil windings, solid insulation
(e.g.\ polyimide, Nomex, epoxy-glass) must be verified for partial-discharge-free
operation at the working dV/dt, which can be severe with hard-switched square waves.

\section{Worked Example: \SI{800}{V}/\SI{800}{V}, \SI{10}{kHz}, \SI{250}{kW}}

\subsection*{Assumptions}
\begin{itemize}
    \item Square-wave excitation (PSFB/DAB-type bridge), $K_f = 4$.
    \item Turns ratio $n=1$ (\SI{800}{V}/\SI{800}{V}).
    \item Nanocrystalline tape-wound core, $B_m = \SI{1.0}{T}$ (derated from
    $B_{sat}\approx\SI{1.2}{T}$ for MF loss margin).
    \item $K_u = 0.3$ (foil winding, generous insulation/creepage at this power level).
    \item $J = \SI{300}{A/cm^2}$ (forced-air or liquid-cooled busbar/foil winding).
\end{itemize}

\subsection*{Step 1 — Apparent power}
\begin{equation}
    P_t = 2P = 2 \times \SI{250}{kW} = \SI{500}{kW}
\end{equation}

\subsection*{Step 2 — Area product}
\begin{equation}
    A_p = \frac{500{,}000 \times 10^4}{4 \times 0.3 \times 1.0 \times 10{,}000 \times 300}
    \approx \SI{1389}{cm^4}
\end{equation}
This is a large core (or, in practice, a paralleled/stacked assembly of standard
nanocrystalline C-cores or cut cores), consistent with a \SI{250}{kW} MF transformer.

\subsection*{Step 3 — Core area and turns}
Assuming a roughly square window/core split, $A_c \approx \sqrt{A_p} \approx \SI{37}{cm^2}$.
From Eq.~\eqref{eq:faraday}:
\begin{equation}
    N_1 = \frac{V_1}{K_f \, A_c \, B_m \, f}
    = \frac{800}{4 \times 37\times10^{-4} \times 1.0 \times 10{,}000}
    \approx 5.4 \;\Rightarrow\; N_1 = N_2 = 6 \ \text{turns}
\end{equation}
(rounding up gives margin against saturation; with $n=1$, $N_2=N_1$).

\subsection*{Step 4 — Conductor sizing}
RMS current per winding (lossless approximation):
\begin{equation}
    I_1 = I_2 = \frac{P}{V} = \frac{250{,}000}{800} = \SI{312.5}{A}
\end{equation}
Required copper cross-section at $J=\SI{300}{A/cm^2}$:
\begin{equation}
    A_{Cu} = \frac{I}{J} = \frac{312.5}{300} \approx \SI{1.04}{cm^2}
\end{equation}
With skin depth $\delta \approx \SI{0.66}{mm}$ at \SI{10}{kHz}, this area should be
realized as a foil (or paralleled foils) of thickness $\le \SI{1.3}{mm}$ and
sufficient width, or an equivalent litz bundle — foil/busbar is the practical
choice at this current level and turn count.

\subsection*{Summary Table}
\begin{center}
\begin{tabular}{@{}ll@{}}
\toprule
Parameter & Value \\
\midrule
Core material & Nanocrystalline (e.g.\ Vitroperm 500F) \\
$B_m$ & \SI{1.0}{T} \\
Area product $A_p$ & $\approx \SI{1389}{cm^4}$ \\
$N_1 = N_2$ & 6 turns \\
$I_1 = I_2$ (RMS) & \SI{312.5}{A} \\
Winding type & Interleaved copper foil, $t \le \SI{1.3}{mm}$ \\
\bottomrule
\end{tabular}
\end{center}

\subsection*{Notes}
This first-pass sizing must be followed by: (i) FEA or analytical leakage
inductance estimation (relevant for DAB/PSFB soft-switching range), (ii) thermal
verification of core and winding losses against the chosen cooling method, and
(iii) insulation coordination for the working voltage and dV/dt. The area-product
method gives the correct order of magnitude but final core selection should match
an actual manufacturer part (or stack thereof) with $A_p \ge$ the calculated value.

\end{document}
